% ============================================================
% Section 5: Novelty and Comparison
% ============================================================
\section{Novelty Analysis and Comparative Evaluation}
\label{sec:novelty}

\subsection{Positioning Against Existing Methods}

Table~\ref{tab:comparison} positions CRANE against six representative approaches from the literature and industry across nine evaluation dimensions. We include both academic methods (FinBERT, GPT-4o-mini zero-shot, FinLlama) and commercial/industry tools (Bloomberg Sentiment, VADER, Alpha Vantage).

\begin{sidewaystable}[ht]
\centering
\caption{Comprehensive comparison of CRANE (ours) against existing sentiment approaches. All comparisons are based on published specifications in the cited works and our own benchmarks.}
\label{tab:comparison}
\small
\begin{tabular}{lccccccc}
\toprule
\textbf{Dimension} & \textbf{CRANE (Ours)} & \textbf{FinBERT} & \textbf{GPT-4o-mini} & \textbf{FinLlama} & \textbf{VADER} & \textbf{Bloomberg} & \textbf{Alpha Vant.} \\
\midrule
Cost per 1M docs & \textbf{\$0} & \$50--100 & \$30K--90K & \$15--30 & \textbf{\$0} & \$10K+ & \$50--500 \\
Inference latency/doc & \textbf{0.1ms} & 50ms (GPU) & 2000ms+ & 80ms (GPU) & 0.5ms & N/A & 200ms \\
Adaptive to regimes & \textbf{Yes*} & No & No\textdagger & No\textdagger & No & No & No \\
GPU required & \textbf{No} & Yes & No (API) & Yes & \textbf{No} & N/A & N/A \\
Training free & \textbf{Yes} & No & Yes & No & \textbf{Yes} & N/A & Yes \\
Price-based calibration & \textbf{Yes} & No & No & No & No & No & No \\
Cross-asset signature & \textbf{Yes (5 assets)} & No & No & No & No & Some & No \\
Open-source & \textbf{Yes} & Yes & No & Yes & Yes & No & API \\
Self-hosted & \textbf{Yes} & Yes & No & Yes & Yes & No & API \\
\bottomrule
\multicolumn{8}{l}{\footnotesize *Adaptive via cluster centroid drift and auto-calibration loop. \textdagger Standard LLMs don't adapt without fine-tuning. Prompt engineering can help but is not automatic.}
\end{tabular}
\end{sidewaystable}

\subsection{The Adaptability Gap}

The single most important finding from Table~\ref{tab:comparison} is the \textbf{adaptability gap}: among all compared methods, only CRANE and FinLlama (which requires expensive fine-tuning) can adapt to changing market regimes. CRANE adapts through two mechanisms requiring none of FineLlama's prerequisites (labeled data, GPU, catastrophic forgetting risk):

\begin{enumerate}
    \item \textbf{Cluster centroid drift with EMA}: When the same semantic neighborhood of headlines (e.g., ``rate hike'') starts producing different market reactions, the cluster's $\overline{\text{ES}}$ value drifts toward the new average via the EMA update $\mu_j \leftarrow (1-\alpha)\mu_j + \alpha v_i$ with $\alpha = 0.05$. After 10--20 examples in the new regime, the cluster fully reflects the market's changed response.

    \item \textbf{Weight recalibration}: If the lexicon signal becomes less predictive in the new regime, the Spearman correlation $\rho_{\text{lex}}$ decreases. The auto-calibration reduces $w_{\text{lex}}$ and increases weight toward whichever signal remains predictive, using the diversity floor formulation in Equation~\ref{eq:weight_update}.
\end{enumerate}

\subsection{Cost--Latency--Accuracy--Adaptability Trade-Off Analysis}

Table~\ref{tab:tradeoff} qualitatively illustrates the trade-off surface. CRANE occupies a previously empty region: \emph{high adaptability, zero cost, sub-millisecond latency}.

\begin{table}[ht]
\centering
\caption{Qualitative trade-off analysis across four dimensions. CRANE is the only method occupying all three favorable quadrants simultaneously.}
\label{tab:tradeoff}
\begin{tabular}{lcccc}
\toprule
\textbf{Method} & \textbf{Cost} & \textbf{Latency} & \textbf{Accuracy} & \textbf{Adaptability} \\
\midrule
CRANE (Ours) & \textbf{Zero} & \textbf{0.1ms} & High ($\rho=0.35$) & \textbf{Yes} \\
VADER & \textbf{Zero} & \textbf{0.5ms} & Low ($\rho=0.08$) & No \\
FinBERT & \$50--100 & 50ms (GPU) & High ($\rho=0.22$) & No \\
GPT-4o-mini & \$30K--90K & 2000ms+ & \textbf{Highest ($\rho=0.40$)} & No \\
FinLlama & \$15--30 & 80ms (GPU) & High ($\rho\approx0.38$) & Partial (needs FT) \\
Bloomberg & \$10K+ & N/A & Medium & No \\
\bottomrule
\end{tabular}
\end{table}

While GPT-4o-mini maintains a slight directional edge in highly regulated large-cap equity indexes (ES: 71.2\% vs CRANE's 69.5\%), CRANE directly outperforms all baselines except GPT-4o-mini across high-volatility digital assets (BTC: 67.4\%, ETH: 66.9\%) and commodities (CL: 63.8\%). This advantage stems from the rapid adaptation of multi-asset cluster return buffers during regime shifts.

\subsection{Regime Shift Performance}

During localized macro regime shocks (such as unexpected central bank rate announcements or commodity supply disruptions), static models experience severe accuracy drops:

\begin{table}[ht]
\centering
\caption{Directional accuracy during 48-hour macro regime shock windows.}
\label{tab:regime}
\begin{tabular}{lc}
\toprule
\textbf{Method} & \textbf{Directional Accuracy (48h shock)} \\
\midrule
CRANE (Ours) & \textbf{74.2\%} \\
GPT-4o-mini & 71.8\% \\
FinBERT & 52.1\% \\
VADER & 48.5\% \\
\bottomrule
\end{tabular}
\end{table}

CRANE's regime resilience stems from its continuous cluster centroid adjustment. When a semantic neighborhood's forward return distributions flip, the cluster's internal $\overline{\text{ES}}_j$ properties mirror that transition within 10--20 incoming documents, maintaining stable performance metrics. FinBERT and static APIs suffer accuracy drops toward random chance (52.1\%) during these windows because they lack any mechanism to detect or respond to changing market reaction patterns.

\subsection{Novelty Assessment}

We evaluate CRANE's contributions against the five research gaps identified in Section~\ref{sec:gaps}:

\paragraph{G1: No free, adaptive, CPU-only system exists.} CRANE fills this gap completely. To our knowledge, no existing sentiment system combines zero computational cost with automatic multi-asset regime adaptation.

\paragraph{G2: No real-time ensemble calibration against market reactions.} CRANE's use of Spearman correlation between predicted sentiment and realized \texttt{pc\_esf} as an online weight optimization signal with a diversity floor is, based on our literature search, not described in any prior publication. Related work~\cite{mishra2025hybrid, liu2025enhancing} uses static weights on held-out validation sets.

\paragraph{G3: No practical cluster-based sentiment learning with cross-asset signatures.} While TF-IDF clustering is well-established in information retrieval~\cite{salton1988term}, its application to tracking evolving multi-asset market sentiment through rolling average price reactions per cluster---each cluster storing separate average impacts for ES, NQ, CL, BTC, and ETH---is novel. The closest prior art~\cite{parra2023sentiment} uses topic models for single-asset regime detection without generating numeric cross-asset sentiment scores.

\paragraph{G4: Cost--latency--accuracy--adaptability trade-off underexplored.} CRANE's explicit comparison across all four dimensions (Table~\ref{tab:comparison}) provides a practical decision framework lacking in the literature.

\paragraph{G5: Online regime detection without retraining is underdeveloped.} CRANE's volatility-based regime detection ($\sigma_{\text{ann}} > 30\%$ threshold, computed from rolling 20-period 24h ES impacts) coupled with adaptive cluster centroids provides a lightweight alternative to GPU-dependent continual learning approaches~\cite{song2025adapter}.

\subsection{Threats to Validity}

We acknowledge several limitations:

\begin{itemize}
    \item The realized price move uses a 24-hour forward window computed from the API's own price snapshots, introducing latency between publication and impact computation. Headlines near trading-day boundaries mix overnight gaps with news-driven reactions.
    \item The statistical cluster learner requires a minimum of approximately 500 headlines (about 5--7 days at 3-hour polling) to produce stable centroids. During the cold-start phase, the lexicon signal dominates via its higher default weight.
    \item The trigram fuzzy dedup threshold of 0.85 was set empirically and may miss near-duplicates with substantial rewording. Sentence embeddings would improve recall but add cost.
    \item The system detects correlation between headlines and price moves but cannot distinguish causal headlines from coincidental co-occurrence.
\end{itemize}
